\chapter{Conclusion}
\label{ch:conclusion}

\section{Conclusion and Future Work}

This capstone investigated liveness verification as a layered engineering problem,
not as a single binary-classification benchmark. The implemented work spans a
lightweight temporal MobileNetV3 mobile path, explicit temporal modelling with
$G^2V^2$former, domain-generalised spatial analysis with GD-FAS, public-data
preparation and transfer experiments, a separate server prototype, and later
manipulation-specific Xception routing. Together these components reveal both the
promise and the current limits of robust remote face verification.

The objectives were met unevenly but informatively. A 24-frame local inference
flow is integrated in Flutter, although its checkpoint provenance and device
performance need a final audit. $G^2V^2$former and GD-FAS were evaluated across
physical-PAD and deepfake domains. Their results show that excellent source-domain
performance can collapse under a change in attack family, and that target-domain
fine-tuning improves some results without solving unseen-domain generalisation.
The later FF++ experiment shows that matched specialists and learned routing can
improve an in-domain universal Xception baseline: the current CNN-routed artefact
raises balanced accuracy from 76.44\% to 84.51\% and ROC-AUC from 86.32\% to
91.68\%. That finding is promising but narrow, because the saved router includes
an original class and has not been tested against unseen generators, injection,
or deployment populations.

The strongest project conclusion is therefore methodological. Physical
presentation attacks, digital manipulation, and capture-channel injection require
different but coordinated evidence. Temporal reasoning can contribute, domain
generalisation must be measured rather than assumed, and manipulation specialists
may help when their routing targets and evaluation are leakage-safe. None of
these replaces genuine-sensor assurance, protected transport, calibrated policy,
privacy protection, subgroup evaluation, or accountable recourse.

The next work should proceed in the following order:

\begin{enumerate}
  \item freeze an experiment manifest and reconcile the MobileNetV3 sequence
        length/checkpoint and $G^2V^2$former metric discrepancies;
  \item rebuild the router as a five-expert selector with no original class,
        using out-of-fold expert predictions if the target is ``best expert'';
  \item compare method-label routing, true best-expert routing, oracle routing,
        top-$k$/soft routing, and a stronger universal or multi-task baseline;
  \item evaluate all deepfake systems at video level on unseen generators,
        additional compression/recapture, and independent external datasets;
  \item calibrate scores and produce reliability, threshold-cost, and escalation-
        coverage curves rather than adopting a universal 0.5 threshold;
  \item connect the mobile and server components through an authenticated,
        versioned contract and implement failure-safe offline/network behaviour;
  \item add genuine-camera confidence, replay protection, device integrity, and
        injection red-team scenarios rather than relying on media analysis alone;
  \item conduct representative device, demographic, accessibility, privacy, and
        reviewer studies with explicit stop criteria and appeal measurement; and
  \item evaluate latency, memory, bandwidth, energy, serving/review cost, licence
        compliance, monitoring, rollback, and secure deletion in a governed pilot.
\end{enumerate}

A human-review path should be implemented only with trained reviewers, bounded
evidence access, reason codes, second-review rules, workload testing, and user
appeal. It should not be described as a generic cure for model uncertainty.

\section{Continuous Engagement and Self-directed Learning}

The project required continuous learning across research and engineering
boundaries. The team moved from an initial mobile/server concept to video
transformers and landmark motion, then to domain-generalisation methods, mobile
model export, API prototyping, Xception specialists, and routing ablations. Reading
primary papers was necessary but insufficient: implementations had to be adapted,
dataset structures reconciled, numerical issues debugged, checkpoints interpreted,
and results compared under common metrics.

Negative transfer results were particularly valuable. They changed the project
from a search for one high score into an investigation of attack taxonomy,
dataset shift, and evidence boundaries. The injection threat similarly expanded
the design beyond classifier accuracy to capture provenance and channel security.
The routing discrepancy reinforced a further lesson: a mathematically attractive
claim is valid only when labels, data splits, code paths, and evaluation actually
implement it.

Professional development should continue through reproducible experiment
packaging, biometric security and privacy study, mobile profiling, secure
deployment, calibration, human-factors evaluation, and responsible disclosure.
Future engagement with a financial partner should begin with jointly defined
threats, lawful and consented data governance, operational error costs, and pilot
stop conditions. These practices are as important to trustworthy liveness
verification as further model optimisation.

